% 第5章图：团{u,v,w}被某个袋容纳的论证示意。T_v={R,A,D,H}, T_u={A,C,D}, T_w={D,G,H}；r_v=R, r_u=A, r_w=D。
\begin{tikzpicture}[x=1mm,y=1mm,>=Stealth,line cap=round,line join=round,
  n/.style={circle,draw=FigureStroke,fill=white,line width=.7pt,minimum size=4.6mm,inner sep=0pt},
  top/.style={n,line width=1.2pt},
  lab/.style={font=\figurefont\footnotesize,text=FigureInk},
  note/.style={font=\FigureNoteFont\fontsize{7.5}{9}\selectfont,text=FigureMuted}]
\coordinate (R) at (0,36); \coordinate (A) at (-20,24); \coordinate (B) at (20,24);
\coordinate (C) at (-32,12); \coordinate (D) at (-8,12); \coordinate (E) at (12,12); \coordinate (F) at (28,12);
\coordinate (G) at (-16,0); \coordinate (H) at (0,0);
% hull strokes: T_v (green), T_u (blue), T_w (yellow)
\draw[FigureGreen,line width=13.2mm] (R)--(A)--(D)--(H);
\draw[FigureGreenFill,line width=12.4mm] (R)--(A)--(D)--(H);
\draw[FigureBlue,line width=9.4mm] (C)--(A)--(D);
\draw[FigureBlueFill,line width=8.6mm] (C)--(A)--(D);
\draw[FigureYellow,line width=6.2mm] (G)--(D)--(H);
\draw[FigureYellowFill,line width=5.4mm] (G)--(D)--(H);
% tree edges
\draw[FigureStroke,line width=.7pt] (R)--(A) (R)--(B) (A)--(C) (A)--(D) (B)--(E) (B)--(F) (D)--(G) (D)--(H);
\foreach \p in {R,A,B,C,E,F,G,H}{\node[n] (n\p) at (\p) {};}
\node[n,line width=1.6pt,minimum size=5.4mm] (nD) at (D) {};
\node[lab,anchor=south] at ($(R)+(0,6.6)$) {根};
\node[lab,anchor=west] at ($(R)+(6.8,0)$) {$r_v$};
\node[lab,anchor=east] at ($(A)+(-5.2,1.5)$) {$r_u$};
\node[lab,anchor=west] at ($(D)+(4.6,1.8)$) {$r_w$};
% legend
\begin{scope}[xshift=44mm,yshift=30mm]
 \fill[FigureGreenFill,draw=FigureGreen,line width=.6pt,rounded corners=1.5pt] (0,0) rectangle (6,4); \node[lab,anchor=west] at (7.5,2) {$T_v$};
 \fill[FigureBlueFill,draw=FigureBlue,line width=.6pt,rounded corners=1.5pt] (0,-8) rectangle (6,-4); \node[lab,anchor=west] at (7.5,-6) {$T_u$};
 \fill[FigureYellowFill,draw=FigureYellow,line width=.6pt,rounded corners=1.5pt] (0,-16) rectangle (6,-12); \node[lab,anchor=west] at (7.5,-14) {$T_w$};
 \node[note,align=left,anchor=north west] at (0,-21) {$r_w$最深，\\$T_u,T_v$都经过$r_w$};
\end{scope}
\end{tikzpicture}
